\vspace{0.5em}\noindent\textbf{Tuần 1 - Phân tích project và kiến trúc cơ
sở}\par

\vspace{0.5em}\noindent\textbf{Mục tiêu}\par

Trong nhóm 5 thành viên, mục tiêu cá nhân của tôi trong tuần đầu là tìm
hiểu các phần mã nguồn liên quan đến triển khai, xác định ranh giới và
phụ thuộc giữa các service, đồng thời lập tài liệu kiến trúc cơ sở cho
công việc Cloud/DevOps ở các tuần sau.

\vspace{0.5em}\noindent\textbf{Phần việc cá nhân}\par

{
\begingroup
\small
\setlength{\tabcolsep}{3pt}
\renewcommand{\arraystretch}{1.15}
\begin{longtable}{@{}
  >{\raggedright\arraybackslash}p{(\linewidth - 4\tabcolsep) * \real{0.3333}}
  >{\raggedright\arraybackslash}p{(\linewidth - 4\tabcolsep) * \real{0.3333}}
  >{\raggedright\arraybackslash}p{(\linewidth - 4\tabcolsep) * \real{0.3333}}@{}}
\toprule\noalign{}
\begin{minipage}[b]{\linewidth}\raggedright
Trạng thái
\end{minipage} & \begin{minipage}[b]{\linewidth}\raggedright
Phần việc được phân công
\end{minipage} & \begin{minipage}[b]{\linewidth}\raggedright
Cơ sở minh chứng
\end{minipage} \\
\midrule\noalign{}
\endhead
\bottomrule\noalign{}
\endlastfoot
Hoàn thành & Khảo sát cấu trúc monorepo và xác định các service chạy
chính. & Cây mã nguồn, \texttt{README.\allowbreak{}md}, \texttt{docker-\allowbreak{}compose.\allowbreak{}yml}
và commit \texttt{4e939cf}. \\
Hoàn thành & Lập bản đồ phụ thuộc giữa frontend, FastAPI backend, chat
service, AI service, PostgreSQL, Redis và DynamoDB. & Các thư mục
service, Compose và cấu hình hệ thống. \\
Hoàn thành & Review luồng request ban đầu giữa browser, API, chat và các
lớp lưu trữ. & API client, backend router, Socket.IO server và storage
service. \\
Hoàn thành & Kiểm tra mô hình khởi chạy local và tài liệu hóa cách chạy
riêng từng service. & Commit \texttt{d3c1168} và \texttt{e43cc04}. \\
\end{longtable}
\endgroup
}

\vspace{0.5em}\noindent\textbf{Triển khai kỹ
thuật}\par

Kho mã nguồn là một monorepo, trong đó mỗi dịch vụ đảm nhiệm một trách
nhiệm riêng:

\begin{itemize}
\tightlist
\item
  \texttt{frontend/\allowbreak{}}: ứng dụng React và Vite dành cho ứng viên và nhân
  sự.
\item
  \texttt{backend/\allowbreak{}}: API FastAPI, mô hình SQLAlchemy, migration Alembic,
  xác thực, công việc, hồ sơ ứng tuyển, tài liệu và các tuyến điều phối
  AI.
\item
  \texttt{chat-\allowbreak{}service/\allowbreak{}}: Node.js, Express, Socket.IO, bộ điều hợp Redis
  và cơ chế lưu hội thoại trên DynamoDB.
\item
  \texttt{ai\_\allowbreak{}service/\allowbreak{}}: phân tích AI, chuẩn hóa CV/công việc, xếp hạng
  lại và mã bộ điều hợp SageMaker được bổ sung sau đó.
\item
  \texttt{docker-\allowbreak{}compose.\allowbreak{}yml}: các thành phần phụ thuộc cục bộ và quy
  trình khởi chạy nhiều dịch vụ.
\end{itemize}

Kiến trúc cơ sở ban đầu:

\textbf{Mô tả sơ bộ.} Trình duyệt sử dụng frontend React/Vite; frontend
gọi backend FastAPI cho các chức năng ứng tuyển và dịch vụ Socket.IO cho
trò chuyện thời gian thực. Backend dùng PostgreSQL, kho tài liệu và dịch
vụ AI, còn chat dùng Redis pub/sub và các bảng DynamoDB.

Kiến trúc cơ sở cũng xác định hướng chuyển đổi lên đám mây ban đầu: đóng
gói backend và dịch vụ chat chạy dài hạn trong container, chuyển dữ liệu
quan hệ bền vững sang PostgreSQL/RDS, dùng DynamoDB cho bản ghi hội
thoại, dùng Redis để phân phối sự kiện thời gian thực, đồng thời chuẩn
bị ứng dụng cho container và CI/CD.

\vspace{0.5em}\noindent\textbf{Vấn đề và giải
pháp}\par

{
\begingroup
\small
\setlength{\tabcolsep}{3pt}
\renewcommand{\arraystretch}{1.15}
\begin{longtable}{@{}
  >{\raggedright\arraybackslash}p{(\linewidth - 6\tabcolsep) * \real{0.2500}}
  >{\raggedright\arraybackslash}p{(\linewidth - 6\tabcolsep) * \real{0.2500}}
  >{\raggedright\arraybackslash}p{(\linewidth - 6\tabcolsep) * \real{0.2500}}
  >{\raggedright\arraybackslash}p{(\linewidth - 6\tabcolsep) * \real{0.2500}}@{}}
\toprule\noalign{}
\begin{minipage}[b]{\linewidth}\raggedright
Vấn đề
\end{minipage} & \begin{minipage}[b]{\linewidth}\raggedright
Nguyên nhân gốc
\end{minipage} & \begin{minipage}[b]{\linewidth}\raggedright
Giải pháp
\end{minipage} & \begin{minipage}[b]{\linewidth}\raggedright
Trạng thái
\end{minipage} \\
\midrule\noalign{}
\endhead
\bottomrule\noalign{}
\endlastfoot
Ảnh chụp mã nguồn đầu tiên chứa phạm vi mã ứng dụng và thành phần phụ
thuộc được sinh tự động quá rộng. & Commit ban đầu đưa toàn bộ nền tảng
vào một commit lớn. & Các công việc sau đó tách tài liệu, tập lệnh và
trách nhiệm thời gian chạy thành các đường dẫn rõ ràng hơn. & Hoàn
thành \\
Dự án ban đầu còn thành phần cơ sở dữ liệu Node liên quan đến Prisma
trong khi backend dùng SQLAlchemy/PostgreSQL. & Chiến lược cơ sở dữ liệu
chuyển sang FastAPI, SQLAlchemy, Alembic và PostgreSQL. & Commit
\texttt{ab60f7d} loại bỏ thành phần Prisma và ghi lại hướng triển khai
EC2/RDS. & Hoàn thành \\
Cách khởi chạy cục bộ cần hỗ trợ nhiều quy trình phát triển. & Một lệnh
duy nhất thuận tiện khi trình diễn, còn các lệnh terminal riêng dễ gỡ
lỗi hơn. & Các commit \texttt{d3c1168} và \texttt{e43cc04} ghi lại cả
hai cách khởi chạy. & Hoàn thành \\
\end{longtable}
\endgroup
}

\vspace{0.5em}\noindent\textbf{Kết quả kiểm thử, build và triển
khai}\par

{
\begingroup
\small
\setlength{\tabcolsep}{3pt}
\renewcommand{\arraystretch}{1.15}
\begin{longtable}{@{}
  >{\raggedright\arraybackslash}p{(\linewidth - 4\tabcolsep) * \real{0.3333}}
  >{\raggedright\arraybackslash}p{(\linewidth - 4\tabcolsep) * \real{0.3333}}
  >{\raggedright\arraybackslash}p{(\linewidth - 4\tabcolsep) * \real{0.3333}}@{}}
\toprule\noalign{}
\begin{minipage}[b]{\linewidth}\raggedright
Hạng mục
\end{minipage} & \begin{minipage}[b]{\linewidth}\raggedright
Kết quả
\end{minipage} & \begin{minipage}[b]{\linewidth}\raggedright
Minh chứng
\end{minipage} \\
\midrule\noalign{}
\endhead
\bottomrule\noalign{}
\endlastfoot
Rà soát mã nguồn & Hoàn thành & \texttt{git\ show\ -\/\allowbreak{}-stat\ 4e939cf}
cho thấy backend, frontend, Docker, migration, kiểm thử và tài liệu được
thêm trong commit đầu tiên. \\
Khởi chạy cục bộ & Hoàn thành một phần & Lịch sử có tập lệnh khởi chạy
và hướng dẫn README nhưng không tìm thấy nhật ký terminal đã lưu. \\
Kiểm thử & Hoàn thành một phần & Có các tệp kiểm thử trong
\texttt{backend/\allowbreak{}tests}; chưa tìm thấy kết quả kiểm thử của tuần 1. \\
Build & Hoàn thành một phần & Có Dockerfile và các tệp dự án Vite; chưa
tìm thấy nhật ký build của tuần 1. \\
Triển khai & Đã lập kế hoạch & Tuần 1 chỉ xác định hướng triển khai,
chưa đặt mục tiêu triển khai AWS. \\
\end{longtable}
\endgroup
}

\vspace{0.5em}\noindent\textbf{Minh chứng}\par

\vspace{0.5em}\noindent\textbf{Commit và yêu cầu kéo
mã}\par

{
\begingroup
\small
\setlength{\tabcolsep}{3pt}
\renewcommand{\arraystretch}{1.15}
\begin{longtable}{@{}
  >{\raggedright\arraybackslash}p{(\linewidth - 6\tabcolsep) * \real{0.2500}}
  >{\raggedright\arraybackslash}p{(\linewidth - 6\tabcolsep) * \real{0.2500}}
  >{\raggedright\arraybackslash}p{(\linewidth - 6\tabcolsep) * \real{0.2500}}
  >{\raggedright\arraybackslash}p{(\linewidth - 6\tabcolsep) * \real{0.2500}}@{}}
\toprule\noalign{}
\begin{minipage}[b]{\linewidth}\raggedright
Commit
\end{minipage} & \begin{minipage}[b]{\linewidth}\raggedright
Mô tả
\end{minipage} & \begin{minipage}[b]{\linewidth}\raggedright
Minh chứng
\end{minipage} & \begin{minipage}[b]{\linewidth}\raggedright
Yêu cầu kéo mã
\end{minipage} \\
\midrule\noalign{}
\endhead
\bottomrule\noalign{}
\endlastfoot
\texttt{4e939cf} & Tạo nền tảng ứng dụng ban đầu gồm backend, frontend,
Docker Compose, migration và kiểm thử. &
\href{https://github.com/Temp-orgo/AWS-Internship/commit/4e939cf0313e73fd915380987cce1a5a7c9728a0}{Xem
commit} & Minh chứng đang chờ \\
\texttt{5ab2924} & Bổ sung đối sánh CV bằng AI và cấu hình AWS RDS. &
\href{https://github.com/Temp-orgo/AWS-Internship/commit/5ab2924ca69ed90f7ef5fe5a454578bfa888e9dc}{Xem
commit} & Minh chứng đang chờ \\
\texttt{ab60f7d} & Loại bỏ thành phần Prisma và ghi lại hướng triển khai
EC2/RDS. &
\href{https://github.com/Temp-orgo/AWS-Internship/commit/ab60f7d02c6c8fc73384198117625dea5973e24f}{Xem
commit} & Minh chứng đang chờ \\
\texttt{d3c1168} & Bổ sung cách khởi chạy toàn bộ hệ thống bằng một
lệnh. &
\href{https://github.com/Temp-orgo/AWS-Internship/commit/d3c1168cea4521823febf32e1fd9714e928d0ddc}{Xem
commit} & Minh chứng đang chờ \\
\texttt{e43cc04} & Ghi lại các lệnh khởi chạy riêng cho từng dịch vụ. &
\href{https://github.com/Temp-orgo/AWS-Internship/commit/e43cc042d868b1ca5c18524d4735dca178f25043}{Xem
commit} & Minh chứng đang chờ \\
\end{longtable}
\endgroup
}

\vspace{0.5em}\noindent\textbf{Nhật ký kiểm
thử}\par

Minh chứng đang chờ: không tìm thấy nhật ký kiểm thử cục bộ hoặc CI của
tuần 1 trong kho mã nguồn.

\vspace{0.5em}\noindent\textbf{Nhật ký build}\par

Minh chứng đang chờ: không tìm thấy nhật ký build Docker hoặc frontend
của tuần 1 trong kho mã nguồn.

\vspace{0.5em}\noindent\textbf{Nhật ký triển
khai}\par

Không áp dụng cho tuần 1. Giai đoạn này bắt đầu lập kế hoạch triển khai;
việc triển khai lên đám mây được dành cho các tuần sau.

\vspace{0.5em}\noindent\textbf{Kết quả trong tuần}\par

Cuối tuần 1, tôi đã hoàn thành bản đồ service/phụ thuộc phục vụ phạm vi
triển khai và xác định hướng container hóa, triển khai AWS ban đầu. Mã
nguồn ứng dụng và các chức năng nền tảng là kết quả chung của nhóm 5
thành viên.

\vspace{0.5em}\noindent\textbf{Bài học rút ra}\par

Bài học chính là kế hoạch triển khai phụ thuộc vào ranh giới thời gian
chạy rõ ràng. Việc sớm tách API, chat, AI, dữ liệu quan hệ, cơ chế
phát/nhận thời gian thực và lưu trữ tài liệu giúp công việc Docker,
Kubernetes và AWS ở các tuần sau dễ phân tích hơn.

\vspace{0.5em}\noindent\textbf{Kế hoạch tuần tiếp
theo}\par

Tuần tiếp theo tập trung vào tính toàn vẹn dữ liệu và xử lý đồng thời:
ràng buộc cơ sở dữ liệu, xử lý yêu cầu trùng, khóa chống lặp, lệnh quy
trình có phiên bản và các kiểm thử chứng minh thao tác ghi đồng thời trả
về xung đột có kiểm soát thay vì dữ liệu không nhất quán.
